%% UniAgent (ICCCM 2026) -- built on the ACM "acmart" sigconf template.
%% For double-anonymous review, change the documentclass line to
%% \documentclass[sigconf, anonymous, review]{acmart}.
%TC:macro \cite [option:text,text]
%TC:macro \citep [option:text,text]
%TC:macro \citet [option:text,text]
%TC:envir table 0 1
%TC:envir table* 0 1
%TC:envir tabular [ignore] word
%TC:envir displaymath 0 word
%TC:envir math 0 word
%TC:envir comment 0 0
% !TeX root = sample-sigconf.tex
\documentclass[sigconf]{acmart}

\AtBeginDocument{%
  \providecommand\BibTeX{{%
    Bib\TeX}}}

%% Figure assets live in paper/images/
\graphicspath{{images/}}
\usepackage{subcaption}
\usepackage{cuted}

%% Give LaTeX more freedom when placing floats so that figures and tables
%% stay on (or nearer) the page that references them instead of drifting.
\setcounter{topnumber}{3}
\setcounter{bottomnumber}{2}
\setcounter{totalnumber}{4}
\setcounter{dbltopnumber}{3}
\renewcommand{\topfraction}{0.92}
\renewcommand{\bottomfraction}{0.5}
\renewcommand{\dbltopfraction}{0.92}
\renewcommand{\textfraction}{0.07}
\renewcommand{\floatpagefraction}{0.72}
\renewcommand{\dblfloatpagefraction}{0.72}

%% dblfloatfix corrects the placement order of deferred double-column
%% (table*/figure*) floats and enables bottom [b] placement, so a wide
%% table is not stranded on the last page above the references.
\usepackage{dblfloatfix}

%% placeins/\FloatBarrier lets us keep a float inside the section that
%% references it (used to pin the comparison table to the Evaluation
%% section instead of letting it drift past the references).
\usepackage{placeins}

%% Rights management information. Replace the SAMPLE values below with
%% the commands provided when the rights form is completed.
\setcopyright{acmlicensed}
\copyrightyear{2026}
\acmYear{2026}
\acmDOI{XXXXXXX.XXXXXXX}
%% ICCCM 2026 proceedings DOI is assigned by ACM after publication.
\acmConference[ICCCM '26]{The 14th International Conference on Computer and Communications Management (ICCCM 2026)}{July 24--26, 2026}{Tokyo, Japan}
\acmISBN{979-8-4007-2439-8}

\begin{document}

\title[UniAgent: End-to-End Decentralized Agent Framework]{UniAgent: An End-to-End Decentralized Framework for Autonomous Agent Discovery, Delegation, and Payment}

\author{Ryoto Taga}
\email{ryoto.taga12@gmail.com}
\affiliation{%
  \institution{Accenture}
  \city{Tokyo}
  \country{Japan}
}
\author{M. Fahim Ferdous Khan}
\email{khan@iniad.org}
\affiliation{%
  \institution{Faculty of Information Networking for Innovation and Design (INIAD)}
  \institution{Toyo University}
  \city{Tokyo}
  \country{Japan}
}

\renewcommand{\shortauthors}{Taga and Khan}

\begin{abstract}
  Large language model (LLM) agents are increasingly able to act on
  behalf of users, but the infrastructure that lets them find, hire,
  and pay one another across organizational boundaries is still
  fragmented. Open protocols already exist for the individual pieces:
  the Agent2Agent (A2A) protocol standardizes agent-to-agent
  communication, x402 turns the HTTP~402 status code into a
  micropayment channel, and ERC-8004 records agent identity on chain.
  However, these protocols have so far been demonstrated only in
  isolation or as one-to-one proofs of concept, and no end-to-end
  service architecture ties them together so that a single
  natural-language request is autonomously decomposed into agent
  discovery, delegation, and payment. In this paper we present
  \emph{UniAgent}, a decentralized agent marketplace that integrates
  A2A, x402, and ERC-8004 under an LLM-based orchestration layer. A
  three-layer design separates a \emph{marketplace layer} (an on-chain
  ERC-8004 registry kept in sync with an off-chain cache through
  blockchain event webhooks), an \emph{orchestration layer} (a ReAct
  agent that discovers, inspects, and invokes external agents through
  three tools), and a \emph{payment layer} (x402 with delegated
  signing). We describe a working
  prototype deployed on the Base Sepolia testnet and report an
  end-to-end evaluation over 20 flight- and hotel-search tasks. UniAgent
  shows that open communication, on-chain identity, and
  protocol-native micropayments can be combined into a coherent,
  vendor-neutral marketplace in which user tasks are carried out
  across separately deployed A2A services.
\end{abstract}

\begin{CCSXML}
<ccs2012>
 <concept>
  <concept_id>10010147.10010178.10010179</concept_id>
  <concept_desc>Computing methodologies~Multi-agent systems</concept_desc>
  <concept_significance>500</concept_significance>
 </concept>
 <concept>
  <concept_id>10002951.10003260.10003282</concept_id>
  <concept_desc>Information systems~Web services</concept_desc>
  <concept_significance>300</concept_significance>
 </concept>
 <concept>
  <concept_id>10003752.10010124.10010131</concept_id>
  <concept_desc>Software and its engineering~Distributed systems organizing principles</concept_desc>
  <concept_significance>300</concept_significance>
 </concept>
</ccs2012>
\end{CCSXML}

\ccsdesc[500]{Computing methodologies~Multi-agent systems}
\ccsdesc[300]{Information systems~Web services}
\ccsdesc[300]{Software and its engineering~Distributed systems organizing principles}

\keywords{AI agents, agent marketplace, A2A protocol, x402, ERC-8004,
  micropayments, on-chain identity, LLM orchestration, ReAct}

\maketitle

\section{Introduction}
\label{sec:introduction}

\paragraph{Problem setting.}
LLM agents can increasingly plan, call tools, and complete multi-step
tasks for users. A natural next step is cross-organizational delegation:
for example, a travel assistant should be able to discover, hire, and
pay specialized flight and hotel agents operated by different providers.
This requires open mechanisms for communication, identity and discovery,
and payment across trust boundaries.

\paragraph{Gap.}
Relevant building blocks now exist: A2A standardizes agent-to-agent
communication~\cite{a2a2025}, ERC-8004 provides on-chain agent
identity~\cite{erc8004}, and x402 attaches micropayments to HTTP
requests~\cite{x402whitepaper}. Yet these protocols do not by themselves
provide an end-to-end marketplace. A2A leaves discovery to applications;
Vaziry et al.~\cite{vaziry2025} demonstrate protocol feasibility only in
fixed one-to-one interactions; and centralized agent stores such as the
GPT Store~\cite{gptstore} provide smooth experiences only inside
platform-specific registration, discovery, and billing systems.
Moreover, combining these protocols is not straightforward: A2A
deliberately excludes discovery from its scope, so a registry mechanism
must be added externally; ERC-8004 provides identity but defines no
interface for capability-based search; and x402 settles payments
independently of agent identity, requiring an integration point that
binds payment authorization to a verified registry entry. No existing
design addresses these gaps together.

\paragraph{Research question.}
This motivates our research question:
\begin{quote}
\itshape How can open agent communication, on-chain identity, and
protocol-level micropayments be integrated into a cohesive marketplace
architecture with LLM-based orchestration, so that a user's high-level
task is autonomously decomposed into agent discovery, delegation, and
payment across organizational boundaries?
\end{quote}

\paragraph{UniAgent overview.}
We answer this question with \emph{UniAgent}, a decentralized agent
marketplace summarized in Figure~\ref{fig:arch}. UniAgent re-implements
the registration, discovery, and payment functions of a closed platform
using A2A, ERC-8004, and x402, and places an LLM ReAct
orchestrator~\cite{yao2023react} above them to analyze the user's task,
discover candidate agents, inspect their specifications, select
providers, and execute paid calls.

\paragraph{Contributions.}
This paper makes three contributions:
\begin{enumerate}
  \item \textbf{An integrated marketplace architecture} that unifies
    A2A, x402, and ERC-8004 behind three loosely coupled layers:
    marketplace, orchestration, and payment.
  \item \textbf{An autonomous LLM pipeline} in which a ReAct agent
    drives task analysis, discovery, specification inspection, and paid
    execution through \texttt{discover\_agents},
    \texttt{fetch\_agent\_spec}, and
    \texttt{execute\_and\_evaluate\_agent}.
  \item \textbf{A working prototype and end-to-end
    demonstration}\footnote{Source code:
    \url{https://github.com/toryoto/UniAgent}; live demo:
    \url{https://uni-agent-web.vercel.app/}.} on Base Sepolia,
    exercising the full flow over external flight- and hotel-search
    services and reporting end-to-end success rate.
\end{enumerate}

\section{Background}
\label{sec:background}

UniAgent builds on four open technologies.
The \emph{Agent2Agent (A2A)} protocol~\cite{a2a2025} standardizes
inter-agent communication: each agent publishes an Agent Card at a
well-known URL describing its capabilities and endpoint, and clients
exchange tasks over JSON-RPC. A2A deliberately leaves agent
\emph{discovery} outside its scope.
The \emph{x402} protocol~\cite{x402whitepaper} revives the HTTP~402
status code for micropayments: a server responds with payment
requirements, the client signs a stablecoin transfer, and a facilitator
settles the transaction on chain.
\emph{ERC-8004}~\cite{erc8004} provides an on-chain agent identity
registry built on ERC-721, assigning each agent a token whose metadata
is stored off chain; because the registry is a public contract, any
client can enumerate agents without trusting a central catalog.
Finally, the \emph{ReAct} pattern~\cite{yao2023react} interleaves LLM
reasoning and tool calling, which suits orchestration where the set of
agents and prices must be explored at run time.

\section{Related Work}
\label{sec:related}

\paragraph{Agent communication protocols.}
Several protocols now compete to standardize how agents interact.
Ehtesham et al.~\cite{ehtesham2025survey} survey this space and compare
the Model Context Protocol (MCP), which connects a single
agent to tools and data sources, with peer-to-peer protocols such as
A2A and the Agent Network Protocol (ANP). UniAgent adopts A2A for
inter-agent communication because it is designed for cross-organization
delegation. None of these communication protocols specifies discovery
or payment, which is precisely what UniAgent's marketplace and payment
layers add.

\paragraph{Decentralized agent identity and closest prior work.}
Vaziry et al.~\cite{vaziry2025} represent each agent's card as an
individual smart contract and show that A2A, x402, and on-chain identity
can interoperate. UniAgent instead consolidates identities into a single
standardized ERC-8004 registry~\cite{erc8004} rather than per-agent
contracts, and---more importantly---builds a full marketplace and
orchestration layer on top, as detailed below.

Vaziry et al. establish the technical \emph{feasibility} of the
protocol combination through one-to-one agent interactions. UniAgent
treats those protocols as building blocks for a marketplace
architecture: a standardized ERC-8004 registry and synchronized cache
support discovery, while an LLM orchestrator starts from a user's
natural-language task, selects providers, pays them over x402, and
returns an integrated result. Thus, Vaziry et al. demonstrate protocol
feasibility, whereas UniAgent contributes the end-to-end service layer
above those protocols.

\paragraph{LLM multi-agent orchestration.}
AutoGen~\cite{wu2023autogen} shows how multiple LLM-based agents can
coordinate through conversation, tool use, and optional human input
within an application. UniAgent addresses a different layer: it treats
the marketplace itself---open registration, cross-vendor discovery, and
protocol-native payment---as the object of orchestration.

\paragraph{Centralized marketplaces.}
Centralized agent stores such as the OpenAI GPT Store~\cite{gptstore}
and Oracle AI Agent Marketplace~\cite{oracleAgentMarketplace2025} show
the demand for integrated agent marketplaces. Oracle's adoption of A2A
also suggests that agent-interface standards are becoming commercially
relevant. However, in these systems, registration, discovery,
deployment, and billing remain confined to vendor platforms and rely on
platform-managed payment rather than protocol-native micropayments.
Rothschild et al.~\cite{rothschild2025agentic} frame the broader tension as a choice
between closed ``agentic walled gardens'' and an open ``web of agents.''
UniAgent targets the latter by combining open protocols for
communication, identity, discovery, and payment.

\section{System Architecture}
\label{sec:architecture}

UniAgent is organized into three layers, shown together in
Figure~\ref{fig:arch}: a \emph{marketplace layer} that maintains the
catalog of available agents, an \emph{orchestration layer} that turns a
user task into a sequence of agent invocations, and a \emph{payment
layer} that settles each invocation. The layers are deliberately
decoupled---the orchestrator does not need to know how identity is
stored, and the payment layer does not need to know why an agent was
selected---so that each can evolve independently. We describe each layer
in turn and then explain how they combine into an end-to-end flow.

\begin{figure*}[t]
  \centering
  \includegraphics[width=0.9\textwidth]{uniagent-system-arch}
  \caption{Overview of the UniAgent architecture. A user issues a
    natural-language task through the web application; an LLM-based
    agent service discovers candidate agents from an ERC-8004 registry
    (cached in PostgreSQL via Alchemy webhooks), invokes the selected
    external A2A agents, and settles payment over x402 using USDC on
    Base Sepolia.}
  \Description{A block diagram with three groups. On the left, a user
    in a browser connects through Privy to the UniAgent application,
    which contains a web app, an agent service, and a PostgreSQL cache.
    The agent service calls the Anthropic Claude API for reasoning and
    communicates with an external A2A agent. On the right, a Base
    Sepolia group contains the EAS, USDC, AgentStaking, and the
    ERC-8004 AgentIdentityRegistry contracts. Arrows show a chat
    execution flow and a marketplace synchronization flow driven by
    Alchemy webhooks and agent-developer registration.}
  \label{fig:arch}
\end{figure*}

\subsection{Marketplace Layer}
\label{sec:marketplace-layer}

The marketplace layer answers a single question: which agents exist and
what can they do? Its authoritative source is the on-chain ERC-8004
\texttt{AgentIdentityRegistry}. An agent developer deploys an A2A
agent, uploads its Agent Card to IPFS, and registers it in the registry
with the IPFS content identifier as the token's metadata URI, optionally
staking a deposit. Reading the chain on every request would be slow and
costly, so UniAgent maintains an off-chain PostgreSQL cache synchronized
through blockchain event webhooks: registration and staking events
trigger a webhook that fetches the metadata from IPFS and updates the
cache. The marketplace UI reads from this cache, while the chain remains
the single source of truth.

\subsection{Orchestration Layer}
\label{sec:orchestration-layer}

The orchestration layer is a ReAct agent~\cite{yao2023react} built with
LangChain and driven by a commercial LLM. It receives the user's
natural-language task and operates the marketplace through three tools:

\begin{sloppypar}
\begin{itemize}
  \item \texttt{discover\_agents(category, query)} returns a short list
    of candidate agents for a given need, drawn from the marketplace
    cache (Section~\ref{sec:marketplace-layer}).
  \item \texttt{fetch\_agent\_spec(agentId)} retrieves the candidate's
    Agent Card from its well-known URL and, when available, its OpenAPI
    specification. This lookup is free of charge and lets the LLM
    inspect skills, input modes, endpoint details, and x402 pricing
    before committing to a paid call.
  \item \texttt{execute\_and\_evaluate\_agent(agentId, category,
    input)} invokes the chosen agent over A2A, settles payment over
    x402 (Section~\ref{sec:payment-layer}), verifies the resulting
    on-chain transaction, and records a quality attestation through EAS.
\end{itemize}
\end{sloppypar}

\paragraph{Orchestration loop.}
Given a task such as planning a trip, the ReAct agent first decomposes
it into sub-tasks (e.g., a flight search and a hotel search). For each
sub-task it calls \texttt{discover\_agents} to obtain up to three ranked
candidates, then typically calls \texttt{fetch\_agent\_spec} on the most
promising ones to compare capabilities and prices. After selecting an
agent, it calls \texttt{execute\_and\_evaluate\_agent}, observes the
result, and repeats the loop until the overall task is complete, at which
point it composes a final answer. Crucially, the model is never given a fixed
plan or a hard-coded list of agents---both the decomposition and the
choice of providers emerge from the ReAct loop at run time.

\paragraph{Discovery ranking.}
The ranking logic is an implementation detail rather than a contribution
of this paper. In the prototype, \texttt{discover\_agents} returns up to
three candidates from the marketplace cache, prioritizing agents with
higher observed reliability and quality, larger deposits, and recent
registration. This keeps the LLM's choice set small while preserving a
basic signal of provider trustworthiness.

\subsection{Payment Layer}
\label{sec:payment-layer}

The payment layer settles each paid invocation using the x402
verify-and-settle pattern~\cite{x402whitepaper}. When the orchestrator
calls an external A2A agent, the agent responds with HTTP~402 and a
\texttt{PAYMENT\allowbreak-\allowbreak REQUIRED} header specifying the price and recipient. The
agent service obtains a delegated payment signature from a Privy wallet
without exposing the user's private key, and retries the request with
the signed payload. The external agent verifies and settles the payment
through an x402 facilitator, which submits the USDC transfer on chain.
The main UniAgent extension to the standard x402 pattern is
\emph{delegated signing} through Privy, which makes payments gasless and
non-custodial. After a successful execution, the system records an
off-chain quality and reliability attestation using the Ethereum
Attestation Service (EAS)~\cite{eas}; these attestations can be used as
reputation signals in future discovery and ranking.

\subsection{Integration Design}
\label{sec:integration-design}

After the user authenticates with Privy and submits a task, the web
application opens a server-sent events (SSE) stream and forwards the
task to the orchestrator. The orchestrator runs the ReAct
loop---discovering agents (marketplace layer), inspecting their Agent
Cards, and executing the chosen agent with x402 payment (payment
layer). Intermediate events---reasoning steps, tool calls, payments,
and partial results---are streamed back to the UI over SSE.

The value of the design comes from the combination rather than any
single protocol. A2A makes agents callable across organizations;
ERC-8004 makes them discoverable without a central directory; x402
makes them payable with a single signed request; and the orchestration
layer ties these together so that a user expresses a goal in natural
language and the system carries out discovery, delegation, and payment
automatically.

\section{Implementation}
\label{sec:implementation}

We implemented UniAgent as a TypeScript monorepo; the source code is
available at \url{https://github.com/toryoto/UniAgent} and the live
deployment at \url{https://uni-agent-web.vercel.app/}. The web application
uses Next.js and React, with Privy for social login and delegated wallets.
The orchestration service runs separately on LangChain with a commercial
LLM, exposing the three tools described in
Section~\ref{sec:orchestration-layer} and streaming results over SSE.
The marketplace cache is stored in PostgreSQL and accessed through an
ORM, with all database access confined to a dedicated data-access layer.
Smart contracts are written in Solidity and developed with Hardhat.
Payments use the x402 protocol with USDC on the Base Sepolia testnet.

\paragraph{Deployment.}
The prototype runs on the Base Sepolia testnet (chain ID~84532), with an
ERC-8004 \texttt{AgentIdentityRegistry} for agent identity and an
\texttt{AgentStaking} contract for deposits. An EAS schema records
per-interaction quality and reliability attestations.

\paragraph{Demonstration scenario.}
For the prototype we focus on a travel-planning use case with two agent
categories, \emph{flight} and \emph{hotel}, served as external A2A
services. A request such as ``plan a three-day trip to Paris'' is
decomposed by the orchestrator into a flight sub-task and a hotel
sub-task, each resolved through discovery, specification inspection, and
paid execution, and the results are combined into a single itinerary.
The evaluated marketplace is intentionally small but heterogeneous. In
each category we deploy one A2A agent backed by a real travel-search API
and ten dummy A2A agents to exercise discovery, ranking, and
provider-quality variation. Thus the evaluation measures protocol-level
cross-service integration rather than production-scale provider coverage.

The web interface includes a chat view for natural-language task
submission and a marketplace view for browsing registered agents by
category and price. During execution, the chat view streams tool calls,
candidate recommendations, human-in-the-loop approval prompts, payment
events, and final results over SSE.

\section{Evaluation}
\label{sec:evaluation}

Our evaluation asks whether UniAgent can carry a realistic user task
through discovery, delegation, and payment end to end. We report 20
completed runs---ten single-agent hotel tasks (S1) and ten multi-agent
flight-plus-hotel tasks (S2)---on Base Sepolia with flight and hotel A2A
services (one API-backed agent and ten dummy agents per category). Each
scenario class uses ten natural-language prompts with varied destinations,
dates, and travel intent.

\subsection{Experimental Setup}
All runs use the same orchestrator prompt, which defines discovery,
specification lookup, user approval, paid execution, and final reporting
without encoding fixed answers or success labels. Before paid execution,
the chat view presents candidate providers and requires explicit user
approval; we therefore measure an approval-mediated end-to-end workflow
rather than autonomous provider selection alone. A run succeeds only if it
completes discovery, specification lookup, candidate presentation, user
approval, paid execution, x402 settlement, and final reporting. We record
transaction hashes and execution logs, then manually label protocol
completion and task success separately.

\subsection{Scenarios and End-to-End Success Rate}
We use two scenario classes. S1 tests a single paid agent call using
hotel-search tasks. A representative prompt is: ``Search for a hotel in
Tokyo for 2 adults from July 24 to July 26, 2026.'' A successful S1 run
requires that one hotel agent be discovered, approved, paid, executed,
and returned in the final response with hotel results.

S2 tests multi-agent composition. The orchestrator must decompose a
travel-planning request into flight and hotel sub-tasks, discover and
inspect providers for each role, execute both providers, and integrate
their results. A representative prompt is: ``Plan a business trip from
Tokyo to London from September 3 to September 5, 2026, including
flights and a business hotel.''

Table~\ref{tab:scenarios} reports the observed completion rates for
these scenario classes.

\begin{table}[H]
  \centering
  \caption{Observed end-to-end results by scenario.}
  \label{tab:scenarios}
  \small
  \begin{tabular}{@{}llccc@{}}
    \toprule
    Scenario & Type & Runs & Protocol & Task \\
    \midrule
    S1 & single agent & 10 & 10/10 & 10/10 \\
    S2 & multi-agent  & 10 & 10/10 & 8/10 \\
    \bottomrule
  \end{tabular}
\end{table}

\paragraph{Representative case.}
S2-02 illustrates the full multi-provider path. For a Tokyo--London
business-trip request, the orchestrator decomposes the task into flight
and hotel searches, discovers and inspects providers for each role,
presents priced candidates for approval, and then executes two paid A2A
calls. A successful run requires two x402 transaction hashes, two agent
results, and a final response that integrates the flight and hotel
options.

The table separates protocol completion from task-level
success. Protocol completion means that UniAgent reached paid execution,
settled the x402 payment, returned transaction hashes, and produced a
final report. Task success additionally requires the external provider
results to match the requested city, dates, and role. All S1 and S2 runs
completed the protocol path. S1 also achieved 10/10 task success. S2
achieved 8/10 task success: in S2-03, the hotel provider returned mainly
Johor Bahru results for a Singapore family trip, and in S2-07 the flight
provider returned a Tokyo--London route for a Tokyo--New York request.
In both cases, UniAgent still completed discovery, payment, execution,
and final reporting, and the orchestrator surfaced the mismatch rather
than presenting the output as fully correct. Minor provider-quality
issues, such as duplicate flight listings or mixed nearby-city hotel
results, were recorded separately from task failures.

Table~\ref{tab:comparison} compares UniAgent feature-by-feature with the
protocol-level study of Vaziry et al.~\cite{vaziry2025} and with
centralized platforms such as the GPT Store~\cite{gptstore} and Oracle
AI Agent Marketplace~\cite{oracleAgentMarketplace2025}.

\begin{table*}[t]
  \centering
  \caption{Feature comparison.}
  \label{tab:comparison}
  \small
  \begin{tabular}{@{}p{0.16\textwidth}p{0.17\textwidth}p{0.17\textwidth}p{0.18\textwidth}p{0.20\textwidth}@{}}
    \toprule
    Feature & Vaziry et al. & UniAgent & GPT Store & Oracle Marketplace \\
    \midrule
    Registration & Custom contract & ERC-8004 & Platform-specific & Validated templates \\
    Discovery & Proposed, not integrated & Registry + cache & In-platform search & Within Fusion Apps \\
    Payment & x402 PoC & x402 + delegated signing & Platform billing & Platform billing \\
    Orchestration & None & LLM ReAct & Platform-dependent & App-embedded \\
    End-to-end flow & No & Yes & Yes (closed) & Yes (closed) \\
    Vendor-neutral & Yes & Yes & No & No \\
    \bottomrule
  \end{tabular}
\end{table*}

\subsection{Comparison with Existing Approaches}
Oracle's adoption of A2A illustrates that agent-interface standards are becoming
commercially relevant, but the benefit remains limited if registration,
discovery, billing, and deployment stay inside a single vendor's
application layer. UniAgent matches the end-to-end
convenience of centralized platforms while keeping registration,
discovery, and payment open and vendor-neutral, and it adds the
orchestration layer that the protocol-level study lacks. The centralized
platforms achieve an end-to-end flow only within closed environments,
whereas UniAgent does so across separately deployed A2A services.

The evaluation results should be read in this context. UniAgent's
strongest measured result is protocol completion: every observed S1 and
S2 run reached paid execution and returned verifiable x402 transaction
hashes. The task-level failures in S2 were not caused by the marketplace,
payment, or A2A integration layers, but by provider output quality. This
distinction is important for an open marketplace architecture. The
framework can make discovery, payment, execution, and reporting
auditable across provider boundaries, while still exposing quality
differences among separately deployed services. In contrast, centralized
marketplaces can hide those boundaries behind a platform-specific
runtime, but they also bind discovery, billing, and orchestration to a
single vendor.

\section{Discussion}
\label{sec:discussion}

The prototype shows that an open, integrated marketplace is feasible:
open registration, discovery, communication, and payment protocols can
be composed into an end-to-end marketplace when connected by an LLM
orchestrator. Its main limitation is scale: the system currently runs on
a testnet, uses two travel-related agent categories, and evaluates a
small population of external A2A services. The discovery ranking and
EAS-based reputation feedback are implemented as design elements but
have not yet been validated under large-scale or adversarial marketplace
conditions.

End-to-end quality also depends on provider output quality and the
orchestrator's selection accuracy. In our evaluation, protocol
completion remained high even when task-level success was reduced by
provider mismatches, suggesting that future work should focus on failure
recovery, stronger reputation signals, mainnet hardening, and larger
multi-agent workflows.

\section{Conclusion}
\label{sec:conclusion}

We presented UniAgent, a decentralized agent marketplace that integrates
the A2A communication protocol, x402 micropayments, and ERC-8004
on-chain identity under an LLM-based orchestration layer. The main
contribution of our architecture is the separation and integration of
three responsibilities that are usually handled by closed platforms: a
marketplace layer for open agent registration and discovery, an
orchestration layer that turns a user request into agent selection and
delegation, and a payment layer that settles cross-provider execution
through protocol-native micropayments. Whereas prior work demonstrated
these protocols in isolation or as one-to-one interactions, UniAgent
shows how they can be composed into a coherent, vendor-neutral framework
in which a single natural-language request is autonomously decomposed
into agent discovery, delegation, and payment. A working prototype on
Base Sepolia demonstrates the full flow over separately deployed flight-
and hotel-search services, showing that end-to-end convenience need not depend
on a centralized agent platform. We believe this integrated design is a
step toward an open ``web of agents'' in which users delegate tasks
across organizational boundaries as easily as they browse the web today.

\bibliographystyle{ACM-Reference-Format}
\bibliography{references}

\end{document}
